\documentclass[14pt]{extarticle}
\usepackage{fontspec}
\usepackage{polyglossia}
\usepackage{amsmath}
\usepackage{amssymb}
\usepackage{xcolor}
\usepackage{fancyhdr}
\usepackage{graphicx}
\usepackage{listings}
\usepackage[most]{tcolorbox}
% \usepackage{setspace}
\usepackage{pifont}
\usepackage{enumitem}
\usepackage{titlesec}
\usepackage[bottom]{footmisc}
\usepackage{titling}
\usepackage{minted}
\usepackage{etoolbox}
\usepackage{array}
\usepackage{longtable} % page-breakable tables (must load before polyglossia/bidi RTL setup)
\usepackage{extsizes}
\usepackage{float}
\usepackage{booktabs}
\usepackage{hyperref}
\usepackage[onehalfspacing]{setspace}
\newif\ifwhiteblack
\ifdefined\setwhiteblack
  \setwhiteblack
\fi

\ifwhiteblack
  \usepackage[a4paper, landscape, top=1.25cm, bottom=1.25cm, left=1.25cm, right=1.25cm, headheight=1.5cm, headsep=0.5cm]{geometry}
\else
  \usepackage[a4paper, top=2.7cm, bottom=2.7cm, left=2.5cm, right=2.5cm, headheight=1.5cm, headsep=0.5cm]{geometry}
\fi
\usepackage{tikz}
\usepackage{datetime2}

\usetikzlibrary{automata, positioning, arrows.meta, arrows}
\providecommand{\LTcaptype}{table}

\newfontfamily\emoji{Segoe UI Emoji}

\pagestyle{fancy}

\setmainlanguage[numerals=western]{arabic}
\setotherlanguage{english}
\newfontfamily\arabicfont[Script=Arabic]{Calibri}
\newfontfamily\arabicfonttt[Script=Arabic]{Courier New}



\ifwhiteblack
\lstset{
  language=[Sharp]C,
  numbers=left,
  stepnumber=1,
  numbersep=8pt,
  frame=single,
  basicstyle=\ttfamily\small,
  keywordstyle=\bfseries,
  stringstyle=\ttfamily,
  commentstyle=\itshape
}
\else
\lstset{
  language=[Sharp]C,
  numbers=left,
  stepnumber=1,
  numbersep=8pt,
  frame=single,
  basicstyle=\ttfamily\small,
  keywordstyle=\color{blue},
  stringstyle=\color{red},
  commentstyle=\color{green!50!black}
}
\fi

\newif\ifdetailed
\ifdefined\setdetailed
  \setdetailed
\fi

\newif\ifwithsols
\ifdefined\setwithsols
  \setwithsols
\fi

\ifwhiteblack
  \AtBeginDocument{
    \LARGE
    \usemintedstyle{bw}
  }
\fi

% Default minted look: compact code that wraps long lines.
% Not applied in black-and-white landscape mode, which already uses \LARGE.
\ifwhiteblack\else
  \setminted{fontsize=\small, breaklines=true}
\fi

% unified tcolorboxes for articles
\ifwhiteblack
\tcbset{colback=white, colframe=black, fonttitle=\bfseries, boxrule=1pt}
\newtcolorbox{boxDef}[1][]{title={{\emoji📘} تعريف\ifx\\#1\\\else ~#1\fi :}}
\newtcolorbox{boxExercise}[1][]{title={{\emoji🧩} تمرين\ifx\\#1\\\else ~#1\fi :}}
\newtcolorbox{boxExample}[1][]{title={{\emoji📝} مثال\ifx\\#1\\\else ~#1\fi :}}
\newtcolorbox{boxNote}[1][]{title={{\emoji✨} ملاحظة\ifx\\#1\\\else ~#1\fi :}}
\newtcolorbox{boxAttention}[1][]{title={{\emoji🔔} تنبيه\ifx\\#1\\\else ~#1\fi :}}
\newtcolorbox{boxWarning}[1][]{title={{\emoji⚡} ملاحظة هامة\ifx\\#1\\\else ~#1\fi :}}
\newtcolorbox{boxSolution}[1][]{title={{\emoji✅} حل\ifx\\#1\\\else ~#1\fi :}}
\newtcolorbox{boxSymbol}[1][]{title={{\emoji🔣} رمز\ifx\\#1\\\else ~#1\fi :}}
\newtcolorbox{boxHint}[1][]{title={{\emoji💡} تلميح\ifx\\#1\\\else ~#1\fi :}}
\else
\tcbset{colback=white, colframe=black, fonttitle=\bfseries, boxrule=0.8pt}
\newtcolorbox{boxDef}[1][]{colback=blue!5!white,colframe=blue!75!black,
  title={{\emoji📘} تعريف\ifx\\#1\\\else ~#1\fi :}}
\newtcolorbox{boxExercise}[1][]{colback=cyan!5!white,colframe=cyan!70!black,
  title={{\emoji🧩} تمرين\ifx\\#1\\\else ~#1\fi :}}
\newtcolorbox{boxExample}[1][]{colback=yellow!5!white,colframe=orange!90!black,
  title={{\emoji📝} مثال\ifx\\#1\\\else ~#1\fi :}}
\newtcolorbox{boxNote}[1][]{colback=gray!10!white,colframe=black,
  title={{\emoji✨} ملاحظة\ifx\\#1\\\else ~#1\fi :}}
\newtcolorbox{boxAttention}[1][]{colback=magenta!10!white,colframe=magenta!80!black,
  title={{\emoji🔔} تنبيه\ifx\\#1\\\else ~#1\fi :}}
\newtcolorbox{boxWarning}[1][]{colback=red!5!white,colframe=red!75!black,
  title={{\emoji⚡} ملاحظة هامة\ifx\\#1\\\else ~#1\fi :}}
\newtcolorbox{boxSolution}[1][]{colback=green!5!white,colframe=green!60!black,
  title={{\emoji✅} حل\ifx\\#1\\\else ~#1\fi :}}
\newtcolorbox{boxSymbol}[1][]{colback=purple!5!white,colframe=purple!70!black,
  title={{\emoji🔣} رمز\ifx\\#1\\\else ~#1\fi :}}
\newtcolorbox{boxHint}[1][]{colback=teal!5!white,colframe=teal!60!black,
  title={{\emoji💡} تلميح\ifx\\#1\\\else ~#1\fi :}}
\fi


\ifwhiteblack
\tcbset{simplecode/.style={ colback=white, colframe=black, boxrule=1pt, arc=2pt, left=4pt,right=4pt,top=4pt,bottom=4pt}}
\else
\tcbset{simplecode/.style={ colback=gray!5, colframe=black!50, boxrule=0.4pt, arc=2pt, left=4pt,right=4pt,top=4pt,bottom=4pt}}
\fi
% Compact code boxes: tighter line spacing (the preamble uses onehalfspacing) and small vertical gaps.
% Must come AFTER the simplecode definition above.
\tcbset{simplecode/.append style={fontupper=\setstretch{0.95}, top=4pt, bottom=4pt, before skip=6pt, after skip=6pt}}
\newenvironment{boxCode}{\begin{tcolorbox}[simplecode]}{\end{tcolorbox}}

% Code cell used inside tables (keeps code left-to-right inside RTL tables).
% Use inside a fixed-width p{...} column: \begin{pycell}\begin{minted}...\end{minted}\end{pycell}
\newtcolorbox{pycell}{frame empty,colback=white,boxsep=0pt,left=2pt,right=2pt,top=2pt,bottom=2pt,width=\linewidth}

\newcolumntype{C}[1]{>{\centering\arraybackslash}p{#1}}

\makeatletter
\DTMnewdatestyle{dotted}{%
  \renewcommand{\DTMdisplaydate}[4]{\two@digits{##3}.\two@digits{##2}.##1}%
  \renewcommand{\DTMDisplaydate}{\DTMdisplaydate}%
}
\makeatother
\DTMsetdatestyle{dotted}

% redefine spaces after titles
\makeatletter
\renewcommand{\@maketitle}{%
  \begin{center}
    {\huge \bfseries \@title \par}%
    \vskip 0.2em % space between title and author
    {\large \@author \par}%
    % \vskip 0.2em % space between author and date
    % {\normalsize \@date \par}%
    {\small \hfill \begin{english}\DTMtoday\end{english} \par}%
  \end{center}
}
\makeatother

\fancyhf{} % clear default
\fancypagestyle{plain}{
  \fancyhf{}
  \fancyhead[L]{مدرسة التسامح الشاملة}
  % \fancyhead[L]{\includegraphics[height=1cm]{../../../images/logoTasamoh.png}}
  \fancyhead[R]{الأستاذ محمود اغبارية}
  \fancyfoot[C]{\thepage}
}

\fancyhead[L]{مدرسة التسامح الشاملة}
\fancyhead[R]{الأستاذ محمود اغبارية}
\fancyfoot[C]{\thepage}

\setcounter{tocdepth}{3} % only section subsection and subsubsection in TOC

\newcommand{\importpdfpage}[2]{%
    \thispagestyle{empty}
    \newgeometry{left=0cm,right=0cm,top=0cm,bottom=0cm}
    \includegraphics[page=#2,width=0.95\paperwidth,height=0.95\paperheight,keepaspectratio]{#1}
    \restoregeometry
}

\newcommand{\insertFullImg}[1]{%
    \noindent
    \makebox[\textwidth][c]{
        \includegraphics[width=0.9\paperwidth,keepaspectratio]{#1}%
    }%
}%


\title{ملخص: حلقات التكرار \textenglish{while} و \textenglish{for} في لغة \textenglish{C\#}}

\begin{document}
\maketitle
\thispagestyle{fancy}

\section*{مقدمة}
الحلقة هي مقطع كود يتكرّر تنفيذه عدّة مرّات لأداء مهمة متكرّرة دون إعادة كتابة الأوامر:
\begin{itemize}[nosep]
    \item التكرار يجعل البرنامج أقصر وأكثر مرونة.
    \item أهم أنواع الحلقات: حلقة \textenglish{while} وحلقة \textenglish{for}.
\end{itemize}

% ==============================================================
% القسم الأول: حلقة while
% ==============================================================
\section{حلقة \textenglish{while}}

\subsection{مفهوم ومبنى حلقة \textenglish{while}}
فكرة الحلقة: \textbf{طالما أنّ الشرط متحقّق (\textenglish{true})، استمرّ في تنفيذ أوامر جسم الحلقة}.
\begin{itemize}[nosep]
    \item نستخدمها عندما لا نعرف عدد التكرارات مسبقًا (التكرار مرتبط بتحقّق شرط معيّن).
    \item يُفحص الشرط \textbf{قبل} كل تكرار؛ فإذا كان \textenglish{false} تتوقف الحلقة فورًا (وإذا كان \textenglish{false} منذ البداية لا يتنفّذ جسم الحلقة ولا مرّة).
\end{itemize}

\begin{boxCode}
\begin{english}
\begin{minted}{csharp}
while (condition)
{
    // Loop body: executed as long as condition is true
}
\end{minted}
\end{english}
\end{boxCode}

% --------------------------------------------------------------
\filbreak
\subsection{استقبال أعداد حتى إدخال قيمة توقّف}
نستقبل القيمة الأولى \textbf{قبل} الحلقة، وفي نهاية جسم الحلقة نستقبل القيمة التالية:

\begin{boxExample}[1: جمع الأعداد المدخلة حتى إدخال $0$]
\begin{english}
\begin{minted}{csharp}
int sum = 0;
int num = int.Parse(Console.ReadLine());
while (num != 0)
{
    sum += num;
    num = int.Parse(Console.ReadLine());
}
Console.WriteLine("Sum = " + sum);
\end{minted}
\end{english}
\end{boxExample}

\filbreak
\begin{boxExample}[2: ضرب الأعداد المدخلة حتى إدخال $1$]
\begin{english}
\begin{minted}{csharp}
int prod = 1;
int num = int.Parse(Console.ReadLine());
while (num != 1)
{
    prod *= num;
    num = int.Parse(Console.ReadLine());
}
Console.WriteLine("Product = " + prod);
\end{minted}
\end{english}
\end{boxExample}

% --------------------------------------------------------------
\filbreak
\subsection{التحقّق من صحة المُدخلات (تكرار الطلب عند الخطأ)}
طالما أنّ المُدخل غير صالح، نطلب من المستخدم إدخال قيمة جديدة:

\begin{boxExample}[إلزام المستخدم بإدخال علامة قانونية بين $0$ و $100$]
\begin{english}
\begin{minted}{csharp}
int grade = int.Parse(Console.ReadLine());
while (grade < 0 || grade > 100)
{
    Console.WriteLine("Invalid! Enter a grade between 0 and 100:");
    grade = int.Parse(Console.ReadLine());
}
Console.WriteLine("Valid grade = " + grade);
\end{minted}
\end{english}
\end{boxExample}

% --------------------------------------------------------------
\filbreak
\subsection{العدّ والتراكم المشروط في \textenglish{while}}

\begin{boxExample}[1: جمع أعداد مدخلة حتى يتجاوز المجموع $10$]
\begin{english}
\begin{minted}{csharp}
int sum = 0;
while (sum <= 10)
{
    int x = int.Parse(Console.ReadLine());
    sum += x;
}
Console.WriteLine("Total = " + sum);
\end{minted}
\end{english}
\end{boxExample}

\filbreak
\begin{boxExample}[2: عدّ الأعداد الزوجية فقط حتى إدخال $-1$]
\begin{english}
\begin{minted}{csharp}
int evenCount = 0;
int num = int.Parse(Console.ReadLine());
while (num != -1)
{
    if (num % 2 == 0)
    {
        evenCount++;
    }
    num = int.Parse(Console.ReadLine());
}
Console.WriteLine("Even count = " + evenCount);
\end{minted}
\end{english}
\end{boxExample}

\filbreak
\begin{boxExample}[3: الاستمرار بالاستقبال حتى قراءة $3$ أعداد زوجية بالضبط]
\begin{english}
\begin{minted}{csharp}
int count = 0;
while (count < 3)
{
    int number = int.Parse(Console.ReadLine());
    if (number % 2 == 0)
    {
        count++;
    }
}
Console.WriteLine("Done reading 3 even numbers.");
\end{minted}
\end{english}
\end{boxExample}

% --------------------------------------------------------------
\filbreak
\subsection{إيجاد القيمة الكبرى أو الصغرى في \textenglish{while}}
\begin{itemize}[nosep]
    \item لإيجاد \textbf{أكبر} قيمة نهيّئ المتغير بـ \textenglish{int.MinValue} (أو بأوّل عدد مدخل).
    \item لإيجاد \textbf{أصغر} قيمة نهيّئ المتغير بـ \textenglish{int.MaxValue} (أو بأوّل عدد مدخل).
\end{itemize}

\begin{boxExample}[إيجاد أصغر عدد مدخل حتى إدخال $0$]
\begin{english}
\begin{minted}{csharp}
int min = int.MaxValue;
int num = int.Parse(Console.ReadLine());
while (num != 0)
{
    if (num < min)
    {
        min = num;
    }
    num = int.Parse(Console.ReadLine());
}
Console.WriteLine("Min = " + min);
\end{minted}
\end{english}
\end{boxExample}

% --------------------------------------------------------------
\filbreak
\subsection{شرط توقّف مركّب في \textenglish{while}}
عندما تستمرّ الحلقة طالما أنّ شرطين متحقّقان معًا (وتتوقف فور انتهاء أحدهما)، نربط بينهما بالعامل \textenglish{\&\&}:

\begin{boxExample}[الاستمرار باستقبال حتى $10$ أعداد كحدّ أقصى أو وصول المجموع إلى $200$]
\begin{english}
\begin{minted}{csharp}
int sum = 0, i = 0;
while (i < 10 && sum < 200)
{
    sum += int.Parse(Console.ReadLine());
    i++;
}
Console.WriteLine("Sum = " + sum + ", Count = " + i);
\end{minted}
\end{english}
\end{boxExample}

\clearpage
% ==============================================================
% القسم الثاني: حلقة for
% ==============================================================
\section{حلقة \textenglish{for}}

\subsection{متى نستخدم \textenglish{for} وما هو مبناها؟}
نستخدم حلقة \textenglish{for} عندما \textbf{نعرف عدد التكرارات مسبقًا} (أو يحدّده متغير مثل $N$).

\begin{boxCode}
\begin{english}
\begin{minted}{csharp}
for (initialization; condition; update)
{
    // Loop body
}
\end{minted}
\end{english}
\end{boxCode}

\begin{itemize}[nosep]
    \item \textenglish{initialization} (التهيئة): تتنفّذ \textbf{مرّة واحدة فقط} في بداية الحلقة (مثل \textenglish{int i = 1}).
    \item \textenglish{condition} (شرط الاستمرار): يُفحص \textbf{قبل} كل تكرار؛ إذا كان \textenglish{true} يتنفّذ الجسم، وإذا كان \textenglish{false} تتوقف الحلقة.
    \item \textenglish{update} (التحديث): يتنفّذ في نهاية كل تكرار \textbf{بعد} تنفيذ جسم الحلقة (مثل \textenglish{i++}).
\end{itemize}

\filbreak
\begin{boxExample}[تتبّع تنفيذ حلقة تطبع الأعداد من $1$ إلى $3$]
\begin{english}
\begin{minted}{csharp}
for (int i = 1; i <= 3; i++)
{
    Console.WriteLine(i);
}
\end{minted}
\end{english}
\end{boxExample}

% --------------------------------------------------------------
\filbreak
\subsection{التكرار عددًا محدّدًا من المرّات والقفزات في العدّاد}
عدّاد الحلقة موجود دائمًا لضبط عدد التكرارات حتى لو لم نستخدمه داخل جسم الحلقة:

\begin{boxExample}[1: طباعة كلمة \textenglish{"Hello"} خمس مرّات]
\begin{english}
\begin{minted}{csharp}
for (int i = 0; i < 5; i++)
{
    Console.WriteLine("Hello");
}
\end{minted}
\end{english}
\end{boxExample}

\filbreak
\begin{boxExample}[2: طباعة الأعداد الزوجية بين $1$ و $100$ بطريقتين]
\begin{english}
\begin{minted}{csharp}
// Method 1: step by 2
for (int i = 2; i <= 100; i += 2)
{
    Console.WriteLine(i);
}

// Method 2: step by 1 and check condition
for (int i = 1; i <= 100; i++)
{
    if (i % 2 == 0)
    {
        Console.WriteLine(i);
    }
}
\end{minted}
\end{english}
\end{boxExample}

\filbreak
\begin{boxExample}[3: طباعة الأعداد التي تقبل القسمة على $5$ من $0$ حتى $N$ بطريقتين]
\begin{english}
\begin{minted}{csharp}
int n = int.Parse(Console.ReadLine());

// Method 1: using if
for (int i = 0; i <= n; i++)
{
    if (i % 5 == 0)
    {
        Console.WriteLine(i);
    }
}

// Method 2: step by 5
for (int i = 0; i <= n; i += 5)
{
    Console.WriteLine(i);
}
\end{minted}
\end{english}
\end{boxExample}

% --------------------------------------------------------------
\filbreak
\subsection{حساب المجموع، القيمة العظمى، والمتوسط في \textenglish{for}}

\begin{boxExample}[1: حساب مجموع الأعداد من $1$ إلى $100$]
\begin{english}
\begin{minted}{csharp}
int sum = 0;
for (int i = 1; i <= 100; i++)
{
    sum += i;
}
Console.WriteLine("Sum = " + sum);
\end{minted}
\end{english}
\end{boxExample}

\filbreak
\begin{boxExample}[2: إيجاد أكبر عدد من بين $10$ أعداد مدخلة]
\begin{english}
\begin{minted}{csharp}
int max = int.MinValue;
for (int i = 0; i < 10; i++)
{
    int x = int.Parse(Console.ReadLine());
    if (x > max)
    {
        max = x;
    }
}
Console.WriteLine("Max = " + max);
\end{minted}
\end{english}
\end{boxExample}

\filbreak
\begin{boxExample}[3: حساب معدّل $5$ أعداد مدخلة]
\begin{english}
\begin{minted}{csharp}
int sum = 0;
for (int i = 0; i < 5; i++)
{
    int x = int.Parse(Console.ReadLine());
    sum += x;
}
double avg = sum / 5.0;
Console.WriteLine("Average = " + avg);
\end{minted}
\end{english}
\end{boxExample}

\clearpage
% ==============================================================
% القسم الثالث: العلاقة بين for و while والحلقات اللانهائية
% ==============================================================
\section{العلاقة بين \textenglish{for} و \textenglish{while} والحلقات اللانهائية}

\subsection{تحويل حلقة \textenglish{for} إلى \textenglish{while}}
كل حلقة \textenglish{for} يمكن كتابتها كحلقة \textenglish{while} مكافئة لها (التهيئة قبل الحلقة، الشرط داخل قوسي \textenglish{while}، والتحديث في نهاية جسم الحلقة):

\begin{boxCode}
\begin{english}
\begin{minted}{csharp}
// Using for:
for (int i = 1; i <= 5; i++)
{
    Console.WriteLine(i);
}

// Equivalent using while:
int i = 1;
while (i <= 5)
{
    Console.WriteLine(i);
    i++;
}
\end{minted}
\end{english}
\end{boxCode}

\filbreak
\subsection{الحلقات اللانهائية وأخطاء شائعة}
تحدث \textbf{الحلقة اللانهائية} عندما يبقى شرط الحلقة \textenglish{true} دائمًا ولا يتغيّر إلى \textenglish{false} أبدًا:

\begin{boxExample}[أمثلة لحلقات لانهائية خاطئة منطقيًّا]
\begin{english}
\begin{minted}{csharp}
// Error 1: variable x is never updated inside while
int x = 5;
while (x > 0)
{
    Console.WriteLine(x);
}

// Error 2: condition (i > 0) stays true forever because i increases
for (int i = 1; i > 0; i++)
{
    Console.WriteLine(i);
}
\end{minted}
\end{english}
\end{boxExample}

\begin{boxWarning}[- تنبيهات هامة]
\begin{itemize}[nosep]
    \item تأكّد دائمًا من تحديث المتغير المؤثّر في شرط التوقف داخل جسم الحلقة (مثل \textenglish{i++} أو قراءة القيمة التالية \textenglish{Console.ReadLine()}).
    \item لا تضع فاصلة منقوطة \textenglish{;} بعد سطر \textenglish{while (...)} أو \textenglish{for (...)} مباشرة.
\end{itemize}
\end{boxWarning}

\vspace{1cm}
\begin{flushleft}
أرجو لكم وقتًا ممتعًا.

الأستاذ محمود اغبارية.
\end{flushleft}

\end{document}
